\subsection{Interval support for a current group: a precise positive construction}
\label{dev:current-support}

\paragraph{The question and its category.}
On P000291, item 1969-180.2, Gelfand asks for an irreducible unitary
representation of an $SU(2)$ current group whose support is an interval
(problem GP-1969-02). The displayed matrices are functions of
$t\in[0,1]$. The notes do not specify their regularity, the topology on
the group, or the continuity required of a representation. Support is
defined through the implication that a current equal to the identity
near a closed set acts trivially. These omissions prevent a single
category-independent answer. They are not silently repaired here.

We give a complete construction for
\[
 G^1=C^1([0,1],SU(2)),\qquad
 G^\infty=C^\infty([0,1],SU(2)),
\]
with their usual $C^1$ and smooth topologies, respectively, and
\emph{strongly continuous} unitary representations. For either group,
say that a closed set $F$ carries a representation $\pi$ if
\[
 g=1\text{ on a neighbourhood of }F\quad\Longrightarrow\quad\pi(g)=1.
\]
Our representation has no proper closed carrying set, so its support in
this precise sense is $[0,1]$. Its restriction to $G^\infty$ remains
irreducible. It is not strongly continuous for the uniform topology.
This is a bounded contribution to the explicitly reconstructed question,
not a classification of all representations, a claim of historical
priority, or an identification of an unstated topology in the notebook.

\paragraph{What the checked literature establishes.}
For operator-norm-continuous irreducible representations of the smooth
current group, the finite-evaluation classification of
\citet[arXiv version, Thms.~4.11 and 4.15]{janssensneeb2015norm}
gives finite support. The preprint circulated on 11 February 2013;
online publication was in December 2014 and the journal issue in 2015.
The group theorem is numbered 4.16 in the published version.
Ismagilov's energy-representation irreducibility theorem requires base
dimension at least five \citep[\S\S1--2]{ismagilov1976unitary}; his
later segment note proves reducibility in dimension one
\citep[pp.~73--74 in the Russian original]{ismagilov1981segment}.
Neither statement rules out the strongly continuous representation
constructed below. Positive-energy localization theorems impose
additional projective and spacetime-action hypotheses
\citep[Introduction]{janssensneeb2024positive}.

These sources document why the continuity and representation class
matter. They do not document why this particular seminar question was
abandoned, nor justify saying that all strongly continuous interval
representations have now been classified. The search included the
original Math-Net notes, the norm-continuous classification and its
references, and the 2024 localization monograph. No historical claim
that the following construction is new is made.

\begin{lemma}[Uniform smooth interpolation at finitely many points]\ours{}
\label{lem:interpolation}
Let $s_1,\ldots,s_q$ be distinct points of $(0,1)$ and set
\[
 r=\frac13\min\bigl(\{s_i,1-s_i:1\leq i\leq q\}
       \cup\{|s_i-s_j|:i\ne j\}\bigr),\qquad L=12/r.
\]
For any $h_1,\ldots,h_q\in SU(2)$ there is a smooth current $g$,
equal to $1$ near the endpoints, with $g(s_i)=h_i$ and
$\|g'\|_\infty\leq L$. The norm is the operator norm on $2\times2$
matrices. If $q=1$, omit the empty set of pairwise distances.
\end{lemma}

\begin{proof}
Every $h_i$ is $\exp X_i$ for a skew-Hermitian traceless matrix $X_i$
of norm at most $\pi$, by unitary diagonalization of an $SU(2)$ matrix.
Use the smooth bump
\[
 \eta(u)=\begin{cases}
 \exp\bigl(1-(1-u^2)^{-1}\bigr),&|u|<1,\\
 0,&|u|\geq1.
 \end{cases}
\]
It satisfies $\eta(0)=1$ and $\|\eta'\|_\infty<3$: putting
$v=(1-u^2)^{-1}\geq1$, its derivative has absolute value at most
$2v^2e^{1-v}\leq8/e<3$. The intervals $(s_i-r,s_i+r)$ are disjoint
and contained in $(0,1)$. Thus
\[
 g(t)=\prod_{i=1}^q\exp\bigl(\eta((t-s_i)/r)X_i\bigr)
\]
has at most one nonidentity factor at any point. It has the required
values, is smooth also at the bump boundaries, and
$\|g'\|_\infty\leq3\pi/r<12/r$.
\end{proof}

\begin{lemma}[Dense pairs with an interpolation tail bound]\ours{}
\label{lem:pairs}
There are pairwise distinct points $a_n,b_n\in(0,1)$, with
$\{a_n:n\geq1\}$ dense, such that, writing
$\epsilon_n=|a_n-b_n|$ and letting $L_N$ be the preceding lemma's
bound for $a_1,b_1,\ldots,a_N,b_N$,
\[
 \sum_n\epsilon_n<\infty,
 \qquad L_N\sum_{n>N}\epsilon_n\leq 2^{-2N}.
 \tag{1}\label{eq:pairtail}
\]
\end{lemma}

\begin{proof}
Enumerate a countable basis $J_1,J_2,\ldots$ of nonempty open
intervals in $(0,1)$. Choose $a_1,b_1\in J_1$ distinct. Inductively
choose $a_m\in J_m$ outside the preceding finite set, and then
$b_m\in J_m$, also new, so close to $a_m$ that
\[
 0<\epsilon_m\leq 2^{-m}
       \min_{1\leq j<m}\frac{2^{-j}}{1+L_j}\quad(m\geq2).
 \tag{2}\label{eq:pairchoice}
\]
All the bounds on the right have already been fixed. Such a choice
is possible in an open interval, even if those bounds are very small.
For fixed $N$ and $m>N$, the $j=N$ term gives
\[
 \sum_{m>N}\epsilon_m
 \leq\frac{2^{-N}}{1+L_N}\sum_{m>N}2^{-m}
 =\frac{2^{-2N}}{1+L_N}.
\]
This proves (1) and summability. Meeting every basic interval proves
density. There is no assumption of a positive separation for the
infinite point set; separation is used only at each finite stage.
\end{proof}

\begin{theorem}[An irreducible representation with interval support]\ours{}
\label{thm:interval-support}
For either $G^1$ or $G^\infty$ as specified above, there is a
strongly continuous unitary representation on a separable,
infinite-dimensional Hilbert space that is irreducible and has support
$[0,1]$. One representation of $G^1$ has all these properties also
upon restriction to $G^\infty$.
\end{theorem}

\begin{proof}
\emph{The Hilbert space and convergence.}
Put $H_n=\mathbb C^2\otimes\overline{\mathbb C^2}$ and
\[
 \Omega_n=2^{-1/2}(e_1\otimes\bar e_1+e_2\otimes\bar e_2).
\]
Here the overline denotes the conjugate Hilbert space and conjugate
representation. The vector $\Omega_n$ is fixed by
$h\otimes\bar h$ for every $h\in SU(2)$. Let $H$ be the completion
of the increasing union $H_1\otimes\cdots\otimes H_N$, embedded at
the next stage by $v\mapsto v\otimes\Omega_{N+1}$. Write
$\Omega=\bigotimes_n\Omega_n$. This defines the incomplete tensor
product directly and makes separability and infinite dimension clear.

For $g\in G^1$ set
\[
 U_n(g)=g(a_n)\otimes\overline{g(b_n)},\qquad
 R_N(g)=\bigotimes_{n\leq N}U_n(g)\otimes1.
\]
Diagonal invariance and the fundamental theorem of calculus give
\[
 \|U_n(g)\Omega_n-\Omega_n\|
 \leq\|g(a_n)-g(b_n)\|\leq\|g'\|_\infty\epsilon_n.
 \tag{3}\label{eq:vacuum-bound}
\]
For a vector $v$ with excitations only in its first $M$ factors,
telescoping a tensor product of unit vectors gives, for $P>Q\geq M$,
\[
 \|R_P(g)v-R_Q(g)v\|
 \leq\|v\|\,\|g'\|_\infty\sum_{Q<n\leq P}\epsilon_n.
\]
The same estimate holds for an arbitrary vector of the first $M$
tensor factors, since that head factors from the vacuum tail.
Summability therefore gives a strong limit $\pi(g)$ first on this
dense union and then on all of $H$, because all $R_N(g)$ have norm
one. Finite products satisfy $R_N(gh)=R_N(g)R_N(h)$.
Strong limits of uniformly bounded products preserve this identity.
Applying it also to $g^{-1}$ proves that $\pi(g)$ is unitary and
$\pi$ is a representation.

\emph{Strong continuity.}
For a vector supported in the first $M$ factors,
\[
 \|\pi(g)v-R_M(g)v\|
 \leq\|v\|\,\|g'\|_\infty\sum_{n>M}\epsilon_n.
 \tag{4}\label{eq:headtail}
\]
As $g\to1$ in $C^1$, the finite head tends to the identity and
$\|g'\|_\infty\to0$. Thus $\pi(g)v\to v$ on the dense union;
unitarity extends this to every vector. Continuity at other group
elements follows from the homomorphism property. The inclusion
$G^\infty\to G^1$ is continuous, so the restriction is strongly
continuous as well.

\emph{Irreducibility.}
Fix $k$ and arbitrary, independently chosen $h_{a_n},h_{b_n}\in SU(2)$
for $n\leq k$. Consider the local unitary
\[
 V=\bigotimes_{n\leq k}
       (h_{a_n}\otimes\overline{h_{b_n}})\otimes1.
\]
For every $N\geq k$, Lemma~\ref{lem:interpolation} supplies a smooth
$g_N$ taking these prescribed values at the first $2k$ points, and
the value $1$ at the remaining points through $2N$, with
$\|g_N'\|_\infty\leq L_N$. For $v$ supported in the first $M$
factors and $N\geq\max(k,M)$, (4) and (1) give
\[
 \|\pi(g_N)v-Vv\|
 \leq\|v\|L_N\sum_{n>N}\epsilon_n
 \leq\|v\|2^{-2N}\longrightarrow0.
\]
Hence $V$ is a strong limit of representation operators, already
from the smooth subgroup. Any operator commuting with the
representation commutes with every such $V$.

The complex linear span of $SU(2)$ in $M_2(\mathbb C)$ is all of
$M_2(\mathbb C)$: the identity and $i$ times the three Pauli matrices
belong to $SU(2)$ and span it. Independent choices at $a_n,b_n$
therefore span $B(H_n)$, and independent choices through $k$ span
the full matrix algebra on the first $k$ factors. If $T$ commutes
with all these local algebras, it commutes with the projections
\[
 Q_N=\bigotimes_{n\leq N}|\Omega_n\rangle\langle\Omega_n|\otimes1.
\]
On the dense union these projections converge strongly to the
rank-one projection onto $\mathbb C\Omega$; hence they do so on
$H$. Consequently $T\Omega=c\Omega$. The vectors $A\Omega$ for
local matrices $A$ span a dense subspace, and
$TA\Omega=AT\Omega=cA\Omega$. Thus $T=c1$. This proves
irreducibility for both groups.

\emph{Support.}
Let $F$ be a proper closed subset of $[0,1]$. Its complement
contains an open interval with closure disjoint from $F$ and lying
in $(0,1)$. Density supplies $a_n$ in that interval. Choose a smooth
bump current $g$, equal to $1$ outside the interval, with $g(a_n)$
noncentral. On the $n$th local matrix algebra, conjugation by
$\pi(g)$ is exactly conjugation by $U_n(g)$: this is true for
$R_N(g)$ when $N\geq n$ and passes to the strong limits of the
unitaries and their inverses. The matrix $U_n(g)$ is not scalar.
Indeed, if a tensor product of two invertible matrices is scalar,
conjugating $X\otimes1$ shows that its first matrix commutes with
every $X$, so that matrix is scalar. Here $g(a_n)$ is noncentral.
It follows that $\pi(g)\ne1$, although $g=1$ near $F$. No proper
closed set carries $\pi$; $[0,1]$ does. Notice that this argument
does not try to isolate one point from the infinite dense set.
\end{proof}

\begin{proposition}[The uniform-topology obstruction]\ours{}
\label{prop:not-czero}
The representation in Theorem~\ref{thm:interval-support}, even
restricted to smooth currents, is not strongly continuous for the
$C^0$ topology. In particular it has no strongly continuous
extension to $C([0,1],SU(2))$ with its uniform topology.
\end{proposition}

\begin{proof}
For each $N$ choose a smooth $f_N:[0,1]\to[0,1]$ with
$f_N(a_j)=1$ and $f_N(b_j)=0$ for $j\leq N$. Disjoint bumps around
the finitely many $a_j$, avoiding the finitely many $b_j$, provide
such a function. With $\sigma_3=\operatorname{diag}(1,-1)$, put
\[
 h_N(t)=\exp(iN^{-1/2}f_N(t)\sigma_3).
\]
Then $\|h_N-1\|_\infty\leq N^{-1/2}$. The first $N$ vacuum
coefficients are all $\cos(N^{-1/2})$, since
\[
 \langle(g(a)\otimes\overline{g(b)})\Omega_n,\Omega_n\rangle
 =\tfrac12\operatorname{Tr}(g(a)g(b)^*).
\]
Every remaining coefficient has absolute value at most one.
Taking limits of the finite tensor products yields
\[
 |\langle\pi(h_N)\Omega,\Omega\rangle|
 \leq\cos(N^{-1/2})^N\longrightarrow e^{-1/2}<1.
\]
Strong convergence $\pi(h_N)\Omega\to\Omega$ is impossible.
An extension continuous for the uniform topology would imply just
that convergence on this sequence.
\end{proof}

\begin{corollary}[Exact support in the norm-continuous class]\ours{}
For $G^\infty$, an irreducible operator-norm-continuous unitary
representation has finite support; conversely, every finite subset
occurs as the support of such a representation. An interval of
positive length cannot occur in this class.
\end{corollary}

\begin{proof}
The smooth mapping group is connected: $g(t)\mapsto g(st)$ joins
it to its constant subgroup, and $SU(2)$ is connected.
The cited classification therefore expresses an irreducible
norm-continuous representation as a finite tensor product
$\bigotimes_{x\in S}\rho_x(g(x))$, omitting trivial factors.
The finite set $S$ carries it. If a closed $F$ omits $x\in S$,
choose a smooth current supported in a small neighbourhood of $x$
disjoint from $F$ and from $S\setminus\{x\}$, taking at $x$ a value
on which the nontrivial $\rho_x$ is not the identity. It acts
nontrivially, so $F$ cannot carry the representation. Thus the
support is exactly $S$. Conversely, independent finite
interpolation and irreducibility of each factor show that the
tensor product of nontrivial finite-dimensional irreducible
$SU(2)$ representations at distinct points is irreducible and
has exactly that support. For $S=\varnothing$ use the trivial
one-dimensional representation.
\end{proof}

\paragraph{Independent check and remaining obstruction.}
The proof was checked by a separate reader and by re-deriving its
two potentially delicate steps. The tail estimate is selected
\emph{after} every finite interpolation constant is fixed, so very
small pair separations do not create a circular bound. Irreducibility
uses approximate realization of finite local matrix algebras, not
impossible exact separation of the dense evaluation set. Support
was checked separately through conjugation of a local algebra.
The $C^0$ obstruction supplies an independent diagnostic that the
topology cannot be weakened without changing the conclusion.
No finite computation is substituted for these infinite-dimensional
arguments.

What remains unresolved here is a classification of strongly
continuous irreducible representations, the identification of the
notebook's intended category, and the extent to which arbitrary
paired sequences yield equivalent representations. A precise next
question for this construction is: for distinct dense pairs
$(a_n,b_n)$ with $\sum|a_n-b_n|<\infty$, characterize when the
corresponding representation is irreducible, without imposing
the sufficient interpolation-tail condition (1). A further bounded
run could compare that condition with intrinsic continuity and
commutant criteria in the infinite-tensor-product literature,
and trace the historical priority of paired evaluation constructions.
It should not promise a classification from the present sufficient
criterion. The status of GP-1969-02 is therefore a positive answer
in the explicit $C^1$/strong category and a negative answer in the
smooth/operator-norm category, with the archival ambiguity retained.

\subsection{Four subspaces: an elementary proof of the requested deformation}
\label{dev:four-subspaces}

On P000317, item 1970-198.12 asks for a simple proof that
$S(5,2)\oplus S(1,2)$ deforms to $S(3,2)\oplus S(3,2)$.
The canonical systems are identified in
\citet[\S2]{gelfandponomarev1971quadruples}, received on
21 October 1970 and published in 1971. The original announcement
was read, as were the later normal forms reproduced in
\citet[\S9]{enomotowatatani2004subspaces}. Those checked texts
identify the systems but do not establish the exact requested
deformation in a located passage. The publication match is therefore
partial. The bounded mathematical question is answered below;
this is not a claim that it remained open after the seminar or
that the argument is new.

\paragraph{Precise formulation.}
Work over $\mathbb C$ with ordered quadruples of subspaces and
simultaneous change of basis in the ambient vector space. The
source's positive-defect notation here means
$\sum_i\dim E_i-2\dim P=2$ for each indecomposable summand;
the notebook's earlier displayed definition has the opposite sign.
Define $S(1,2)$ as four copies of $\mathbb C$ in $\mathbb C$.
For $R=S(3,2)$ take ambient basis $e_1,e_2,f_1$ and
\[
 \begin{aligned}
 E_1&=\langle e_1,e_2\rangle,&
 E_2&=\langle f_1,e_2\rangle,\\
 E_3&=\langle e_1,e_2+f_1\rangle,&
 E_4&=\langle f_1,e_1+e_2\rangle.
 \end{aligned}
\]
For $S(5,2)$ use basis $e_1,e_2,e_3,f_1,f_2$ and
\[
 \begin{aligned}
 E_1&=\langle e_1,e_2,e_3\rangle,&
 E_2&=\langle f_1,f_2,e_3\rangle,\\
 E_3&=\langle e_1,e_2+f_1,e_3+f_2\rangle,&
 E_4&=\langle f_1,e_1+f_2,e_2+e_3\rangle.
 \end{aligned}
\]
Both direct sums in the question consequently define points of
$X=\operatorname{Gr}(4,6)^4$. We prove the stronger statement
that $R\oplus R$ is the unique dense open orbit in $X$; thus every
point of $X$, including the displayed special sum, admits an
analytic deformation to that orbit. The field, topology and
ordered-subspace convention are explicit reconstruction choices.

\begin{lemma}[Endomorphisms of the three-dimensional system]\ours{}
Every endomorphism of $R$ preserving its four subspaces is scalar.
\end{lemma}

\begin{proof}
In the dual basis the four annihilator lines are represented by
\[
 \ell_1=(0,0,1),\quad\ell_2=(1,0,0),\quad
 \ell_3=(0,1,-1),\quad\ell_4=(1,-1,0).
\]
The first three form a basis, and
$\ell_4=-\ell_1+\ell_2-\ell_3$ has all three coefficients nonzero.
An endomorphism preserving each primal plane has transpose
preserving these four dual lines. It is diagonal in the first
three lines, and preserving the fourth forces all three diagonal
entries to agree. This proves the assertion, including
noninvertible endomorphisms.
\end{proof}

\begin{theorem}[The notebook's deformation]\ours{}
The $\mathrm{GL}_6(\mathbb C)$ orbit of $R\oplus R$ is open and
dense in $X$. There is a holomorphic one-parameter family of ordered
quadruples $x(t)\in X$, defined near $0$, with
$x(0)=S(5,2)\oplus S(1,2)$ and $x(t)\cong R\oplus R$ for all
sufficiently small nonzero $t$.
\end{theorem}

\begin{proof}
An endomorphism of a direct sum is a block matrix of morphisms
between its summands. The lemma gives
$\operatorname{End}(R\oplus R)=M_2(\mathbb C)$, of dimension four.
Its invertible elements form the stabilizer, also of dimension
four. The orbit therefore has dimension $36-4=32$. Meanwhile
$\dim X=4\cdot4(6-4)=32$, and $X$ is irreducible. An algebraic
group orbit is locally closed. A locally closed subset of full
dimension in this irreducible variety has closure $X$ and is open
in that closure; the orbit is consequently open dense. Two
disjoint open dense orbits in an irreducible variety cannot exist.

For completeness, the closure statement gives an actual
one-parameter deformation, not just a dimension comparison.
Choose a product of standard affine Grassmann charts through
$x(0)$; in graph coordinates it is $\mathbb C^{32}$ with $x(0)=0$.
The complement of the open orbit is a proper closed algebraic
subset of this chart. Choose a nonzero polynomial $f$ vanishing
on that complement and a point $v$ with $f(v)\ne0$.
Then $f(tv)$ is a nonzero one-variable polynomial, so for every
sufficiently small $t\ne0$ it is nonzero. The graph coordinates
$tv$ give the asserted holomorphic family. Finally, the common
intersection of the four subspaces is one-dimensional for
$S(5,2)\oplus S(1,2)$ and zero-dimensional for $R\oplus R$;
the special fibre is genuinely a different orbit.
\end{proof}

\begin{example}[Independent endomorphism calculation]\oursComputed{}
The script \nolinkurl{code/verify_four_subspaces.py} imposes
$\ell^T A v=0$ for every subspace basis vector $v$ and every
annihilating covector $\ell$, then row-reduces the resulting
linear system over $\mathbb Q$. It independently finds
\[
\begin{array}{c|rrrrr}
\text{system}&S(1,2)&S(3,2)&S(5,2)&S(5,2)\oplus S(1,2)&R\oplus R\\
\hline
\dim\operatorname{End}&1&1&1&7&4
\end{array}
\]
The last two orbit dimensions are $29$ and $32$. The output is
\nolinkurl{output/verify_four_subspaces.json}. This calculation
checks the normal forms and stabilizer dimensions independently
of the projective-frame proof. The first two nontrivial single
summands have full-dimensional orbits in their own products of
Grassmannians, confirming the source's stability claim for these
particular summands, without endorsing its separate false slogan
about every parameter-free indecomposable.
\end{example}

The proof was independently reviewed through both the projective
frame and the orbit-dimension argument. No obstruction remains
for this particular deformation over $\mathbb C$. What remains
unverified is the exact historical publication of the requested
elementary proof, not a mathematical openness claim. One further
bounded run could locate that attribution or produce a convenient
closed formula for a chosen one-parameter change of bases;
neither is needed for the existence theorem proved here.

\subsection{Audit annotation: commuting quantum marginals need not have product eigenbases}
\label{dev:bell}

This section corrects item 1970-197.2 on P000312--P000313. It is an
annotation of a false literal claim, not a newly discovered open
problem selected from class C or D. The notes first discuss the
simple-spectrum situation, where a matrix commuting with all
one-site marginals is diagonal in their uniquely determined tensor
eigenbasis, and then assert that a suitable tensor eigenbasis can
always be chosen in the degenerate case. The following counterexample
checks the stronger condition that \emph{all} finite marginals
commute after embedding in their common tensor algebra.

\begin{proposition}[A faithful consistent Bell-block counterexample]\ours{}
For each $0<p<1$ there is a state on the qubit chain indexed by
$\mathbb Z$ whose finite density matrices are strictly positive,
consistent under partial trace, and pairwise commuting when
embedded into a common finite tensor algebra, but whose two-site
density matrix on a designated adjacent pair is not diagonal in
any product of one-site orthonormal bases.
\end{proposition}

\begin{proof}
Partition the sites into pairs $\{2k,2k+1\}$ and put on each pair
\[
 D_p=p|\Phi\rangle\langle\Phi|+\frac{1-p}{4}1_4,
 \qquad \Phi=\frac{|00\rangle+|11\rangle}{\sqrt2}.
\]
Its eigenvalues are $(1+3p)/4$ on $\Phi$ and $(1-p)/4$ on its
three-dimensional orthogonal complement. Thus it is positive
definite and has trace one. Each one-site partial trace is $1_2/2$.
For a finite set $I\subset\mathbb Z$, define $\rho_I$ as the
tensor product of $D_p$ for complete pairs in $I$ and $1_2/2$ for
single sites of partially included pairs. Factors are reordered
into site order by the canonical tensor permutation.

Tracing out one member of a complete pair replaces $D_p$ by
$1_2/2$; tracing out the remaining member gives one. It follows
that $\operatorname{Tr}_{I\setminus J}\rho_I=\rho_J$ for every
$J\subset I$. These normalized positive consistent functionals
define a state on the increasing union of local matrix algebras;
their norm is one, so the state extends to its norm completion.
Every finite density matrix is strictly positive.

For finite $I,J$, embed $\rho_I$ and $\rho_J$ in the tensor
algebra on $I\cup J$ by tensoring with identity operators. On each
fixed two-site block, each embedded matrix is either $D_p$ or a
scalar multiple of the identity (including the case with only one
site in the union). These block factors commute, hence so do the
two full matrices. This checks every finite $I,J$, not just nested
intervals.

If $D_p$ were diagonal in a product of one-site orthonormal bases,
its simple eigenvalue $(1+3p)/4$ would have a product eigenvector.
Its eigenspace is instead $\mathbb C\Phi$. A two-qubit product
vector has a rank-one coefficient matrix, whereas $\Phi$ has
coefficient matrix $2^{-1/2}1_2$ of determinant $1/2$. Thus
$\Phi$ is not a product vector, a contradiction.
\end{proof}

The example has period two. No one-site translation-invariance
hypothesis is stated at this point in the notebook, and none is
needed for the claim being refuted. Strict positivity shows that
zero eigenvalues are not responsible. The obstruction is
degeneracy of the one-site marginals, which leaves room for a
nonproduct eigenvector of a joint density matrix.

\begin{proposition}[The valid simple-spectrum repair]\ours{}
Let $\rho$ be a density matrix on
$H_1\otimes\cdots\otimes H_m$, with finite-dimensional $H_i$.
Suppose that each one-site marginal $\rho_i$ has simple spectrum
and that $\rho$ commutes with each embedded $\rho_i$. Then $\rho$
is diagonal in the tensor product of the one-site eigenbases.
\end{proposition}

\begin{proof}
For a simple-spectrum self-adjoint matrix, each spectral projection
is a polynomial in that matrix by Lagrange interpolation on its
distinct eigenvalues. Thus $\rho$ commutes with every one-site
rank-one spectral projection and with every product of these
projections. Those products are the rank-one projections of the
tensor eigenbasis, so each corresponding line is invariant under
$\rho$. This is the desired diagonalization.
\end{proof}

\begin{example}[Independent exact finite checks]\oursComputed{}
The file \nolinkurl{code/verify_annotations.py} uses rational
arithmetic from the Python standard library. On four sites it
computes every one of the $16$ subset marginals by an explicit
basis definition of partial trace, independently compares them
with the block formula in the proof, and checks all $256$ ordered
commutators after embedding. It repeats the entire calculation
for $p=1/2$ and $p=2/3$. Both runs pass exactly. The first parameter
gives spectrum $5/8,1/8,1/8,1/8$, and the second gives
$3/4,1/12,1/12,1/12$. As an additional diagnostic, the partial
transpose has a negative eigenvalue $-1/8$ and $-1/4$, respectively.
The full output is \nolinkurl{output/verify_annotations.json}.
The infinite-chain conclusion follows from the block proof above;
four-site computations are an independent check of its algebra,
not evidence substituted for an all-volume theorem.
\end{example}

\subsection{Audit annotation: the product uncertainty question and wavelet scaling}
\label{dev:wavelet}

On P000358 the general orthogonal-system question asks for a bound
on the \emph{product} $\delta_x\delta_p$; it is distinct from the
proposed Gram--Schmidt construction from a coherent-state lattice.
The item and GP-1970-59 retain that distinction. This is a class-B
historical resolution, not an open problem manufactured for
development. Meyer's Schwartz orthonormal wavelet construction
\citep[\S2, Thm.~1]{meyer1987uncertainty} already gives the
following elementary implication. The lecture was delivered in
February 1986 and published in 1987.

\begin{proposition}[Uniform uncertainty products from a Schwartz wavelet]\ours{}
Let $\psi\in\mathcal S(\mathbb R)$ have norm one and suppose that
\[
 \psi_{j,k}(x)=2^{j/2}\psi(2^jx-k),\qquad j,k\in\mathbb Z,
\]
is an orthonormal basis of $L^2(\mathbb R)$. With
$Qf=xf$, $Pf=-if'$ and $\hbar=1$, let
$\Delta_f A=\|(A-\langle Af,f\rangle)f\|$ for a normalized vector
in the domain of $A$. Then
\[
 \Delta_{\psi_{j,k}}Q=2^{-j}\Delta_\psi Q,
 \qquad\Delta_{\psi_{j,k}}P=2^j\Delta_\psi P,
 \qquad
 \Delta_{\psi_{j,k}}Q\,\Delta_{\psi_{j,k}}P
 =\Delta_\psi Q\,\Delta_\psi P<\infty.
\]
\end{proposition}

\begin{proof}
All vectors are Schwartz, so the required first and second moments
exist. In the position integrals substitute $y=2^jx-k$ to obtain
\[
 \langle Q\psi_{j,k},\psi_{j,k}\rangle
 =2^{-j}(\langle Q\psi,\psi\rangle+k),
\]
and the centered second moment is multiplied by $2^{-2j}$.
Differentiation gives
$P\psi_{j,k}(x)=2^{3j/2}(P\psi)(2^jx-k)$, so the momentum mean is
multiplied by $2^j$ and its centered second moment by $2^{2j}$.
Taking square roots proves the identities.
\end{proof}

As a second check, the dilation--translation unitary
$U_{j,k}f(x)=2^{j/2}f(2^jx-k)$ satisfies on Schwartz functions
\[
 U_{j,k}^*QU_{j,k}=2^{-j}(Q+k),\qquad
 U_{j,k}^*PU_{j,k}=2^jP.
\]
These conjugation identities reproduce the variance calculation
without repeating the moment integrals. They also show its exact
limitation: neither individual deviation stays bounded as
$j\to\pm\infty$.

Bourgain proved the stronger existence of unrestricted orthonormal
bases with both centered second moments uniformly bounded
\citep[main theorem; IHES preprint, November 1986]{bourgain1988uncertainty}.
The journal paper appeared in 1988. The distinction between such
bases, regular Gabor systems, and higher-moment obstructions is
explicit in \citet[pp.~115--116 and Thm.~1.1]{grochenigmalinnikova2013localization}.
Neither the wavelet argument nor that general existence statement
proves the notebook's particular Gram--Schmidt procedure works in
an unspecified ordering. The chapter accordingly separates the
settled existence question from that incompletely specified method.
